\problemname{Triangelfabriken}
\noindent

Tristian jobbar på en triangelfabrik.
Hans arbete är att klassificera olika sorters trianglar som produceras i fabriken.
Tristian ber nu dig att skriva ett program som kan ersätta honom.

Tristian kommer att ge dig 3 positiva heltal, $a$, $b$ och $c$, som är de 3 vinklarna i en triangel.

Ditt program ska göra följande:
\begin{itemize}
  \item Om triangeln är trubbig, så ska programmet skriva ut ``Trubbig Triangel''.
  \item Om triangeln är spetsig, så ska programmet skriva ut ``Spetsig Triangel''.
  \item Om triangeln är rätvinklig, så ska programmet skriva ut ``Ratvinklig Triangel''.
  (Tristian har inte lärt sig alla bokstäver i det svenska alfabetet än, så han förstår inte om du skriver ut ``Rätvinklig'' med ett ``ä''. Han förstår bara ``Ratvinklig Triangel'').
\end{itemize}

För att påminna om vad trubbiga, spetsiga och rätvinkliga trianglar är:
\begin{itemize}
  \item En triangel är trubbig om någon av dess vinklar är större än $90$ grader.
  \item En triangel är spetsig om alla dess vinklar är mindre än $90$ grader.
  \item En triangel är rätvinklig om en av dess vinklar är exakt $90$ grader.
\end{itemize}

\section*{Indata}
Indatan består av tre rader, där varje rad består av ett heltal.

Den första raden innehåller ett heltal $a$ $(1 \leq a < 180)$, som beskriver vinkeln i antal grader på det första hörnet i triangeln.
Den andra raden innehåller ett heltal $b$ $(1 \leq b < 180)$, som beskriver vinkeln i antal grader på det andra hörnet i triangeln.
Den tredje raden innehåller ett heltal $c$ $(1 \leq c < 180)$, som beskriver vinkeln i antal grader på det tredje hörnet i triangeln.

Det är garanterat att indatan är en triangel, alltså att $a + b + c = 180$.

\section*{Utdata}
Skriv ut vilken sorts triangel det är. Svaret måste antingen vara ``Trubbig Triangel'', ``Spetsig Triangel'', eller ``Ratvinklig Triangel''.

\section*{Poängsättning}
Din lösning kommer att testas på en mängd testfallsgrupper.
För att få poäng för en grupp så måste du klara alla testfall i gruppen.

\noindent
\begin{tabular}{| l | l | p{12cm} |}
  \hline
  \textbf{Grupp} & \textbf{Poäng} & \textbf{Gränser} \\ \hline
  $1$    & $100$  & Inga ytterligare begränsningar. \\ \hline
\end{tabular}

